% 通用作业模板 (中文) — 可直接用于 XeLaTeX / pdfLaTeX

\documentclass[UTF8]{ctexart}
\usepackage[a4paper,margin=2.5cm]{geometry}
\usepackage{amsmath,amssymb}
\usepackage{graphicx}
\usepackage{hyperref}
\usepackage{caption}
% 可选页眉页脚
% \usepackage{fancyhdr}
% \pagestyle{fancy}

% ---- 元信息（请在文档开头设置这些命令） ----
\newcommand{\course}{示例课程}
\newcommand{\assignment}{第*次作业}
\newcommand{\studentname}{路诚钺}
\newcommand{\studentid}{202518019435002}
\newcommand{\due}{\today}

% ---- 文档开始 ----
\begin{document}

\begin{center}
  {\Huge \bfseries \course}\\[1ex]
  {\large \assignment}\\[2ex]
  \rule{0.6\textwidth}{0.6pt}\\[2ex]
  姓名：\underline{\studentname} \quad 学号：\underline{\studentid} \\
  日期：\due
\end{center}

\vspace{1.5em}

% 简短说明：在下面按题号填写答案。
\section*{解答}

% 范例题目与答题区
\subsection*{第1题}
\textbf{题目：} 在此粘贴或写出题目。

\textbf{解：}
\vspace{1em}

% 建议的子结构：推导、计算、结论
\subsubsection*{推导}
在此书写推导过程。可使用数学环境：
\[ E=mc^2 \]

\subsubsection*{计算 / 结果}
在此给出计算过程与最终结果。

\subsubsection*{结论}
在此写结论或答题要点。

% 复用多个题目时复制上面的块。

\vfill
% 可选附录或图表
%\section*{附录}
%\begin{figure}[h]
%  \centering
%  \fbox{\parbox[c][6cm][c]{10cm}{在此放置图像或截图（使用 \texttt{\includegraphics}）}}
%  \caption{示例图}
%\end{figure}

\end{document}
